\chapter{Pochodne wyższych rzędów}
\label{ma1:ch:higher-derivatives}

Pochodna funkcji jest sama funkcją, którą można ponownie różniczkować. W ten sposób otrzymujemy informacje o zmianach prędkości, przyspieszeniu i dalszych zmianach lokalnych. W następnym rozdziale wykorzystamy pochodne wyższych rzędów do przybliżeń wielomianowych.

\section{Definicja i klasy regularności}
\begin{definicja}
\label{ma1:def:higher-derivatives}
Niech $I\subseteq\R$ będzie niepustym przedziałem otwartym, a $f:I\to\R$. Przyjmujemy $f^{(0)}=f$ i rekurencyjnie definiujemy $f^{(n+1)}=(f^{(n)})'$ wszędzie tam, gdzie $f^{(n)}$ jest określona w otoczeniu punktu i różniczkowalna w tym punkcie. Jeśli wszystkie pochodne do rzędu $n$ istnieją na $I$, funkcję nazywamy $n$ razy różniczkowalną na $I$. Piszemy $f\in C^n(I)$, gdy pochodne do rzędu $n$ istnieją i są ciągłe na $I$. W szczególności $C^0(I)=C(I)$ oraz $C^\infty(I)=\bigcap_{n\ge0} C^n(I)$.
\end{definicja}
Jeżeli $f$ jest $n$ razy różniczkowalna na $I$, gdzie $n\ge1$, to $f^{(j)}$ jest różniczkowalna dla $0\le j<n$, a zatem ciągła na mocy twierdzenia \ref{ma1:thm:linear-approximation}. Otrzymujemy $f\in C^{n-1}(I)$, lecz niekoniecznie $f\in C^n(I)$. W szczególności
\[
 C^\infty(I)\subseteq C^{n+1}(I)\subseteq C^n(I).
\]
Na przedziale domkniętym $[a,b]$, $a<b$, zapis $f\in C^n([a,b])$ oznacza tutaj ciągłość $f$ na $[a,b]$, przynależność ograniczenia do $C^n((a,b))$ i istnienie ciągłych przedłużeń wszystkich pochodnych rzędów $1,\ldots,n$ do końców. Z twierdzenia \ref{ma1:thm:derivative-endpoint-limit}, zastosowanego kolejno do ciągłych przedłużeń pochodnych niższych rzędów, wartości tych przedłużeń na końcach są odpowiednimi pochodnymi jednostronnymi.

\section{Liniowość i wzór Leibniza}
\begin{twierdzenie}
\label{ma1:thm:higher-linearity}
Jeśli $f,g:I\to\R$ są $n$ razy różniczkowalne na otwartym przedziale $I$, to dla $\alpha,\beta\in\R$ funkcja $\alpha f+\beta g$ jest $n$ razy różniczkowalna i
$(\alpha f+\beta g)^{(n)}=\alpha f^{(n)}+\beta g^{(n)}$.
\end{twierdzenie}
\begin{proof}
Dla $n=0$ teza jest definicją. Jeśli formuła zachodzi do rzędu $k<n$, to jej prawa strona jest różniczkowalna na $I$, gdyż istnieją $f^{(k+1)}$ i $g^{(k+1)}$. Reguła liniowości pochodnej (twierdzenie \ref{ma1:thm:derivative-algebra}) daje po zróżniczkowaniu formułę rzędu $k+1$. Indukcja kończy dowód.
\end{proof}

\begin{twierdzenie}[Wzór Leibniza]
\label{ma1:thm:leibniz}
Przy tych samych założeniach iloczyn $fg$ jest $n$ razy różniczkowalny i
\[
 (fg)^{(n)}=\sum_{k=0}^n \binom nk f^{(k)}g^{(n-k)}.
\]
\end{twierdzenie}
\begin{proof}
Dla $n=0$ otrzymujemy $fg=fg$. Załóżmy formułę rzędu $m<n$. Każdy składnik jest różniczkowalny, ponieważ występujące w nim pochodne mają rzędy co najwyżej $m$. Z reguły iloczynu oraz liniowości wynika
\[
 (fg)^{(m+1)}
 =\sum_{k=0}^m\binom mk f^{(k+1)}g^{(m-k)}
 +\sum_{k=0}^m\binom mk f^{(k)}g^{(m+1-k)}.
\]
W pierwszej sumie podstawiamy $j=k+1$, a w drugiej $j=k$. Wyrazy skrajne to $fg^{(m+1)}$ oraz $f^{(m+1)}g$. Dla $1\le j\le m$ współczynnik przy $f^{(j)}g^{(m+1-j)}$ wynosi $\binom m{j-1}+\binom mj=\binom{m+1}j$, na mocy tożsamości Pascala (wynikającej z podziału podzbiorów $(m+1)$-elementowego zbioru według obecności wyróżnionego elementu). Oba współczynniki skrajne są równe $1$. Dostajemy formułę rzędu $m+1$ i kończymy indukcję.
\end{proof}
\begin{przyklad}
\label{ma1:ex:leibniz-polynomial-exp}
Dla $n\ge2$ i $f(x)=x^2e^x$ jedynie trzy składniki wzoru Leibniza mogą być niezerowe, ponieważ pochodne $x^2$ od rzędu trzeciego znikają. Zatem
\[
 f^{(n)}(x)=e^x\bigl(x^2+2nx+n(n-1)\bigr).
\]
Dla $n=0,1$ ta sama formuła zachodzi bezpośrednio z definicji i reguły iloczynu.
\end{przyklad}

\section{Potęgi i funkcje elementarne}
\begin{stwierdzenie}
\label{ma1:prop:higher-elementary}
Dla $p\in\{0,1,\ldots\}$ oraz $n\ge0$ mamy
\[
 (x^p)^{(n)}=
 \begin{cases}
 \dfrac{p!}{(p-n)!}x^{p-n},&n\le p,\\
 0,&n>p.
 \end{cases}
\]
Na $(0,\infty)$, dla dowolnego $\alpha\in\R$,
$(x^\alpha)^{(n)}=\alpha(\alpha-1)\cdots(\alpha-n+1)x^{\alpha-n}$, z pustym iloczynem równym $1$. Ponadto dla $a,b\in\R$:
\[
 \begin{aligned}
 (e^{ax+b})^{(n)}&=a^n e^{ax+b},\\
 (\sin(ax+b))^{(n)}&=a^n\sin(ax+b+n\pi/2),\\
 (\cos(ax+b))^{(n)}&=a^n\cos(ax+b+n\pi/2).
 \end{aligned}
\]
Dla $x>0$ i $n\ge1$ zachodzi $(\log x)^{(n)}=(-1)^{n-1}(n-1)!/x^n$. W formułach z $a^n$ przyjmujemy $a^0=1$ również dla $a=0$.
\end{stwierdzenie}
\begin{proof}
Dla potęg naturalnych zaczynamy od rzędu $0$ i stosujemy regułę potęgową z wniosku \ref{ma1:cor:integer-power-derivative}. Jeśli $n<p$, różniczkowanie mnoży współczynnik przez $p-n$, a wykładnik zmniejsza o $1$. W rzędzie $p$ otrzymujemy stałą $p!$, której wszystkie dalsze pochodne znikają. Dla $\alpha\in\R$ i $x>0$ z definicji $x^\alpha=\exp(\alpha\log x)$ oraz reguły łańcuchowej dostajemy $(x^\alpha)'=\alpha x^{\alpha-1}$. Stosując tę równość do wykładnika $\alpha-n$ w każdym kroku, otrzymujemy podany iloczyn; obowiązuje on również wtedy, gdy któryś czynnik jest zerowy.

Dla funkcji wykładniczej i trygonometrycznych przypadek $n=0$ jest definicją. Reguła łańcuchowa oraz wzory na pochodne z rozdziału \ref{ma1:ch:derivative} dają w każdym kroku dodatkowy czynnik $a$. Dla sinusa używamy $\cos t=\sin(t+\pi/2)$, a dla cosinusa $-\sin t=\cos(t+\pi/2)$; przesunięcie argumentu zwiększa się więc o $\pi/2$. Jeśli $a=0$, wszystkie trzy funkcje są stałe, zgodnie ze wzorami dla $n\ge1$. Dla logarytmu przypadek $n=1$ to $1/x$. Różniczkując $(-1)^{n-1}(n-1)!x^{-n}$, otrzymujemy $(-1)^n n!x^{-n-1}$, czyli wzór dla $n+1$.
\end{proof}

\section{Złożenia i odwrotności}
\begin{twierdzenie}
\label{ma1:thm:higher-composition}
Niech $I,J$ będą otwartymi przedziałami, $g:I\to J$, $f:J\to\R$. Jeżeli $f,g$ są $n$ razy różniczkowalne na swoich dziedzinach, to $f\circ g$ jest $n$ razy różniczkowalna. Jeśli ponadto $f,g$ należą do odpowiednich klas $C^n$, to $f\circ g\in C^n(I)$.
\end{twierdzenie}
\begin{proof}
Dla $n=0$ pierwsza teza nie nakłada warunków, a druga wynika z ciągłości złożenia. Załóżmy obie tezy dla $n=m$ i niech $f,g$ będą $(m+1)$ razy różniczkowalne. Reguła łańcuchowa (twierdzenie \ref{ma1:thm:chain-rule}) daje $(f\circ g)'=(f'\circ g)g'$. Funkcje $f'$ i $g$ są co najmniej $m$ razy różniczkowalne, podobnie $g'$. Z założenia indukcyjnego $f'\circ g$ jest $m$ razy różniczkowalna, a wzór Leibniza daje tę własność iloczynu. Stąd istnieją pochodne $f\circ g$ do rzędu $m+1$. Jeśli $f,g\in C^{m+1}$, to $f'\circ g$ oraz $g'$ są klasy $C^m$ na $I$, a ich iloczyn jest tej samej klasy: wszystkie jego pochodne do rzędu $m$ są skończonymi sumami iloczynów ciągłych funkcji. Ponadto $f\circ g$ jest ciągła, jako złożenie funkcji ciągłych. Otrzymujemy $f\circ g\in C^{m+1}(I)$.
\end{proof}
Dwukrotne i trzykrotne zastosowanie reguł łańcuchowej oraz iloczynu daje
\[
 \begin{aligned}
 (f\circ g)''&=(f''\circ g)(g')^2+(f'\circ g)g'',\\
 (f\circ g)'''&=(f'''\circ g)(g')^3
     +3(f''\circ g)g'g''+(f'\circ g)g'''.
 \end{aligned}
\]
Istotnie, różniczkowanie pierwszego składnika wzoru na drugą pochodną daje $(f'''\circ g)(g')^3+2(f''\circ g)g'g''$, a drugiego $(f''\circ g)g'g''+(f'\circ g)g'''$; po zsumowaniu otrzymujemy trzecią pochodną.

\begin{wniosek}
\label{ma1:cor:higher-inverse}
Jeśli $n\ge1$, $f\in C^n(I)$ oraz $f'(x)\ne0$ na otwartym przedziale $I$, to $J=f(I)$ jest otwartym przedziałem, $f$ jest bijekcją $I\to J$ i $h=f^{-1}$ należy do $C^n(J)$. Dla $n\ge2$ zachodzi
\[
 h''(y)=-\frac{f''(h(y))}{(f'(h(y)))^3}.
\]
\end{wniosek}
\begin{proof}
Ciągłość $f'$ oraz twierdzenie o wartości pośredniej wykluczają zmianę znaku bez zera. Z twierdzenia \ref{ma1:thm:derivative-monotonicity} funkcja $f$ jest ściśle monotoniczna. Ciągłość zapewnia, że $J$ jest przedziałem. Dla każdego $x\in I$ można wybrać $u<x<v$ w $I$; ścisła monotoniczność umieszcza $f(x)$ ściśle między $f(u)$ i $f(v)$, więc $J$ zawiera otoczenie $f(x)$. Jest zatem otwarty. Odwrotność $h$ jest ciągła: dla $\varepsilon>0$ wybieramy $u<x<v$ z $x-u,v-x<\varepsilon$, a wartości $y$ dostatecznie bliskie $f(x)$ leżą między $f(u)$ i $f(v)$, skąd $u<h(y)<v$. Twierdzenie \ref{ma1:thm:inverse-derivative} daje $h'=1/(f'\circ h)$. Prawa strona jest ciągła i dobrze określona, więc $h\in C^1(J)$.

Jeżeli $1\le k<n$ i $h\in C^k(J)$, to $f'\in C^k(I)$, a twierdzenie \ref{ma1:thm:higher-composition} daje $f'\circ h\in C^k(J)$. Funkcja odwrotności multiplikatywnej $t\mapsto 1/t$ jest gładka na każdej z półosi (ze wzorów na potęgi całkowite i reguły różniczkowania); argument $f'\circ h$ ma stały niezerowy znak. Kolejne zastosowanie twierdzenia o złożeniach daje $h'\in C^k(J)$, skąd $h\in C^{k+1}(J)$. Indukcja daje tezę. Różniczkując $h'=(f'\circ h)^{-1}$, dostajemy $h''=-(f'\circ h)^{-2}(f''\circ h)h'$, a podstawienie wzoru na $h'$ daje żądaną równość.
\end{proof}

\section{Ścisłość hierarchii regularności}
\begin{przyklad}
\label{ma1:ex:strict-c-hierarchy}
Dla $m\ge0$ funkcja $F_m(x)=x^m|x|$ należy do $C^m(\R)$, ale nie do $C^{m+1}(\R)$. Dla $x>0$ jest równa $x^{m+1}$, a dla $x<0$ jest równa $-x^{m+1}$. Dla $0\le j\le m$ jej pochodne poza zerem mają zatem postać $\pm (m+1)!x^{m+1-j}/(m+1-j)!$, ze znakiem zależnym od półosi. W zerze $F_m(0)=0$. Jeżeli $j<m$ i $F_m^{(j)}(0)=0$, to iloraz różnicowy w zerze jest na obu półosiach równy $\pm (m+1)!h^{m-j}/(m+1-j)!$, a więc dąży do zera. Indukcja daje $F_m^{(j)}(0)=0$ dla $j\le m$. Wszystkie te pochodne są ciągłe w zerze, ponieważ wykładnik $m+1-j$ jest dodatni. Jednak $F_m^{(m)}(x)=(m+1)!|x|$, więc w zerze ilorazy jednostronne wynoszą $\pm(m+1)!$ i pochodna rzędu $m+1$ nie istnieje.
\end{przyklad}

\begin{przyklad}
\label{ma1:ex:second-derivative-discontinuous}
Niech $G(0)=0$ i $G(x)=x^4\sin(1/x)$ dla $x\ne0$. Iloraz $G(h)/h=h^3\sin(1/h)$ dąży do zera, zatem $G'(0)=0$. Poza zerem
$G'(x)=4x^3\sin(1/x)-x^2\cos(1/x)$.
Ponieważ $G'(h)/h=4h^2\sin(1/h)-h\cos(1/h)\to0$, mamy $G''(0)=0$. Dla $x\ne0$ reguły iloczynu i łańcuchowa dają
\[
 G''(x)=12x^2\sin(1/x)-6x\cos(1/x)-\sin(1/x).
\]
Dla $x_k=1/(\pi/2+2k\pi)$ oraz $y_k=1/(3\pi/2+2k\pi)$, $k\ge1$, oba ciągi dążą do zera, a wartości $G''(x_k)$ dążą do $-1$, natomiast $G''(y_k)$ do $1$. Funkcja jest więc dwukrotnie różniczkowalna na $\R$, lecz nie jest klasy $C^2(\R)$.
\end{przyklad}

\section{Interpretacja mechaniczna}
Jeśli $s(t)$ opisuje położenie punktu poruszającego się wzdłuż prostej, to $v=s'$ jest prędkością (ze znakiem), $a=s''$ przyspieszeniem, a $s'''$ szybkością zmian przyspieszenia. Przy stałym przyspieszeniu $a_0$ i warunkach $s(0)=s_0$, $s'(0)=v_0$ na przedziale otwartym zawierającym $0$ mamy $s(t)=s_0+v_0t+a_0t^2/2$. Istotnie, różnica $s'(t)-a_0t$ ma pochodną zerową, więc jest stała z twierdzenia \ref{ma1:thm:derivative-monotonicity}; wartość w zerze daje $v_0$. Następnie różnica $s(t)-v_0t-a_0t^2/2$ także ma zerową pochodną i jest równa swojej wartości $s_0$ w zerze. Nie wymaga to jeszcze teorii całki.

\begin{center}
\begin{tikzpicture}[x=1.3cm,y=1.05cm,>=Stealth]
 \draw[->] (-.15,0)--(2.35,0) node[right] {$t$};
 \draw[->] (0,-.15)--(0,4.5) node[above] {$s$};
 \draw[blue!65!black,thick,domain=0:2,samples=60] plot (\x,{\x*\x});
 \node[blue!65!black,above left] at (1.7,2.89) {$s=t^2$};
 \begin{scope}[xshift=4.3cm]
 \draw[->] (-.15,0)--(2.35,0) node[right] {$t$};
 \draw[->] (0,-.15)--(0,4.5) node[above] {$v$};
 \draw[orange!85!black,thick] (0,0)--(2,4);
 \node[orange!85!black,above left] at (1.7,3.4) {$v=2t$};
 \end{scope}
 \begin{scope}[xshift=8.6cm]
 \draw[->] (-.15,0)--(2.35,0) node[right] {$t$};
 \draw[->] (0,-.15)--(0,4.5) node[above] {$a$};
 \draw[green!40!black,thick] (0,2)--(2,2);
 \node[green!40!black,above] at (1,2) {$a=2$};
 \end{scope}
\end{tikzpicture}

\small Rysunek: położenie, prędkość i przyspieszenie w ruchu jednostajnie przyspieszonym, przy $s_0=v_0=0$ i $a_0=2$; pokazano przedział $0\le t\le2$.
\end{center}
